%% ============================================================================
%% HeatLoad HVAC Calculation Tool — Project Documentation
%% Generated: 2026-06-06
%% ============================================================================

\documentclass[12pt, a4paper, oneside]{report}

% ── Packages ──────────────────────────────────────────────────────────────────
\usepackage[utf8]{inputenc}
\usepackage[T1]{fontenc}
\usepackage{lmodern}
\usepackage[margin=1in]{geometry}
\usepackage{graphicx}
\usepackage{xcolor}
\usepackage{hyperref}
\usepackage{listings}
\usepackage{booktabs}
\usepackage{longtable}
\usepackage{tabularx}
\usepackage{amsmath}
\usepackage{amssymb}
\usepackage{enumitem}
\usepackage{fancyhdr}
\usepackage{tocloft}
\usepackage{titlesec}
\usepackage{caption}
\usepackage{float}
\usepackage{textcomp}
\usepackage{microtype}
\usepackage{parskip}

% ── Color Definitions ────────────────────────────────────────────────────────
\definecolor{codeblue}{RGB}{0,102,204}
\definecolor{codegray}{RGB}{128,128,128}
\definecolor{codegreen}{RGB}{0,128,0}
\definecolor{codered}{RGB}{204,51,0}
\definecolor{backcolour}{RGB}{248,248,248}
\definecolor{accentblue}{RGB}{41,98,255}
\definecolor{warnorange}{RGB}{230,126,34}
\definecolor{sectionblue}{RGB}{26,54,126}

% ── Hyperref Setup ────────────────────────────────────────────────────────────
\hypersetup{
  colorlinks=true,
  linkcolor=accentblue,
  urlcolor=codeblue,
  citecolor=codegreen,
  pdftitle={HeatLoad HVAC Calculation Tool — Documentation},
  pdfauthor={Engineering Team},
  pdfsubject={ASHRAE Heat Load Calculation — React Application},
}

% ── Listings Setup ────────────────────────────────────────────────────────────
\lstdefinestyle{mystyle}{
  backgroundcolor=\color{backcolour},
  commentstyle=\color{codegreen}\itshape,
  keywordstyle=\color{codeblue}\bfseries,
  numberstyle=\tiny\color{codegray},
  stringstyle=\color{codered},
  basicstyle=\ttfamily\footnotesize,
  breakatwhitespace=false,
  breaklines=true,
  captionpos=b,
  keepspaces=true,
  numbers=left,
  numbersep=8pt,
  showspaces=false,
  showstringspaces=false,
  showtabs=false,
  tabsize=2,
  frame=single,
  rulecolor=\color{codegray!30},
  xleftmargin=1.5em,
  framexleftmargin=1em,
}
\lstset{style=mystyle}

% ── Header/Footer ────────────────────────────────────────────────────────────
\pagestyle{fancy}
\fancyhf{}
\fancyhead[L]{\small\textcolor{codegray}{HeatLoad HVAC Tool}}
\fancyhead[R]{\small\textcolor{codegray}{\leftmark}}
\fancyfoot[C]{\thepage}
\renewcommand{\headrulewidth}{0.4pt}
\renewcommand{\footrulewidth}{0pt}

% ── Title Format ──────────────────────────────────────────────────────────────
\titleformat{\chapter}[display]
  {\normalfont\huge\bfseries\color{sectionblue}}{\chaptertitlename\ \thechapter}{20pt}{\Huge}
\titleformat{\section}
  {\normalfont\Large\bfseries\color{sectionblue}}{\thesection}{1em}{}
\titleformat{\subsection}
  {\normalfont\large\bfseries\color{sectionblue!80}}{\thesubsection}{1em}{}

% ── Custom Commands ───────────────────────────────────────────────────────────
\newcommand{\file}[1]{\texttt{\textcolor{codeblue}{#1}}}
\newcommand{\code}[1]{\texttt{\textcolor{codered}{#1}}}
\newcommand{\ashrae}{\textsc{ASHRAE}}
\newcommand{\warn}[1]{\textcolor{warnorange}{\textbf{Warning:}} #1}
\newcommand{\note}[1]{\textcolor{accentblue}{\textbf{Note:}} #1}

% ══════════════════════════════════════════════════════════════════════════════
%                              DOCUMENT BEGIN
% ══════════════════════════════════════════════════════════════════════════════
\begin{document}

% ── Title Page ────────────────────────────────────────────────────────────────
\begin{titlepage}
  \centering
  \vspace*{2cm}
  {\Huge\bfseries\textcolor{sectionblue}{HeatLoad HVAC Calculation Tool}\par}
  \vspace{0.5cm}
  {\Large\textcolor{codegray}{Comprehensive Project Documentation}\par}
  \vspace{1.5cm}
  {\large\textcolor{codegray}{Version 2.8 — June 2026}\par}
  \vspace{2cm}

  \rule{\textwidth}{0.4pt}
  \vspace{0.5cm}

  {\large
  \begin{tabular}{rl}
    \textbf{Framework:}   & React 19 + Redux Toolkit + Vite 7 \\
    \textbf{Language:}     & TypeScript / JavaScript (JSX/TSX) \\
    \textbf{Styling:}      & Tailwind CSS v4 \\
    \textbf{Deployment:}   & Vercel (SPA) \\
    \textbf{Standards:}    & ASHRAE HOF 2021, ASHRAE 62.1-2022, ISO 14644-1:2015 \\
    \textbf{Industries:}   & Semiconductor, Pharmaceutical, Battery, Solar \\
  \end{tabular}
  }

  \vspace{2cm}
  \rule{\textwidth}{0.4pt}
  \vfill
  {\small Generated automatically from source code analysis.}
\end{titlepage}

% ── Table of Contents ─────────────────────────────────────────────────────────
\tableofcontents
\newpage

% ══════════════════════════════════════════════════════════════════════════════
\chapter{Project Overview}
% ══════════════════════════════════════════════════════════════════════════════

\section{Introduction}

The \textbf{HeatLoad HVAC Calculation Tool} is a client-side, single-page application (SPA) designed for HVAC engineers working in critical-environment industries---semiconductor fabrication, pharmaceutical manufacturing, battery manufacturing (Li-ion and lead-acid), and solar/PV production. The application automates the full \ashrae{} Handbook of Fundamentals (2021) Chapter~18 heat load calculation methodology, from building envelope analysis through psychrometric state-point resolution.

The tool replaces traditional Excel-based Room Data Sheet (RDS) calculators with a reactive, browser-based interface that provides instant recalculation as engineers modify room parameters, climate data, or system design assumptions.

\section{Key Engineering Capabilities}

\begin{itemize}[leftmargin=1.5em]
  \item \textbf{Multi-Season Load Calculation} --- Summer, monsoon, and winter seasons with peak-season auto-selection for both CFM (supply air) and TR (cooling capacity).
  \item \textbf{CLTD/CLF Envelope Method} --- Walls, roofs, glazing, skylights, partitions, and floors with ASHRAE construction presets.
  \item \textbf{Psychrometric Engine} --- Hyland-Wexler saturation pressure (ASHRAE HOF 2021 Ch.1, Eq.~3 \& 5) replacing the less-accurate Magnus approximation.
  \item \textbf{ASHRAE 62.1-2022 Ventilation} --- Full Ventilation Rate Procedure (VRP) with industry-specific categories including lead-acid battery rooms (OSHA 29~CFR~1926.403(i)).
  \item \textbf{ISO 14644-1:2015 Classification} --- ISO~1 through ISO~9, CNC, and Unclassified with IEST-RP-CC012.2 ACPH ranges.
  \item \textbf{GMP Annex 1:2022 Compliance} --- Grade~A through D mapping with mandated ACPH floors and 1.25$\times$ process safety factors.
  \item \textbf{Altitude-Corrected Psychrometrics} --- All factors ($C_s$, $C_l$, humidity ratio) corrected via the barometric formula to actual site elevation.
  \item \textbf{Reheater Sizing} --- Physics-based minimum ESHF calculation and ASHRAE Ch.18~\S17.3 reheat load determination.
  \item \textbf{Required ADP Analysis} --- ESHF line / saturation curve intersection with three-outcome classification (\code{found}, \code{sensible\_only}, \code{no\_solution}).
  \item \textbf{CHW/HW Pipe Sizing} --- Automated pipe diameter selection with velocity-based branch and manifold sizing.
  \item \textbf{Humidification Sizing} --- Winter moisture deficit calculation with desiccant system awareness for battery dry rooms.
\end{itemize}

\section{Technology Stack}

\begin{table}[H]
\centering
\caption{Technology Stack Summary}
\begin{tabularx}{\textwidth}{lX}
\toprule
\textbf{Technology} & \textbf{Purpose} \\
\midrule
React 19.2            & UI component framework \\
Redux Toolkit 2.11    & Centralised state management with Immer-powered immutable updates \\
React Router DOM 7.13 & Client-side routing (SPA) \\
Vite 7.3              & Build tooling, HMR dev server \\
Tailwind CSS 4.2      & Utility-first CSS framework \\
TypeScript 6.0        & Static type checking \\
Vitest 4.1            & Unit testing framework \\
ExcelJS 4.4           & Excel export (\texttt{.xlsx}) generation \\
Lucide React          & Icon library \\
React Hot Toast       & Notification toasts \\
Vercel                & Deployment platform with SPA rewrite rules \\
\bottomrule
\end{tabularx}
\end{table}


% ══════════════════════════════════════════════════════════════════════════════
\chapter{Architecture}
% ══════════════════════════════════════════════════════════════════════════════

\section{High-Level Architecture}

The application follows a \textbf{unidirectional data flow} pattern:

\begin{enumerate}
  \item \textbf{User interaction} dispatches Redux actions.
  \item \textbf{Redux reducers} (slices) produce new immutable state.
  \item \textbf{Memoised selectors} (\code{createSelector}) derive computed RDS data.
  \item \textbf{React components} re-render with the new data.
\end{enumerate}

All calculation logic is \textbf{pure}---implemented as selector modules in \file{features/results/}. These modules are never registered as Redux reducers; they derive from the five state slices and produce the full RDS row object on every state change.

\section{State Shape}

The Redux store is composed of five independent slices:

\begin{lstlisting}[language={}, caption={Root Reducer Composition}]
rootReducer = {
  project:  projectReducer,   // Project metadata, ambient, system design
  ahu:      ahuReducer,        // Air Handling Unit definitions
  room:     roomReducer,       // Room list and active selection
  climate:  climateReducer,    // Outdoor design conditions (3 seasons)
  envelope: envelopeReducer,   // Per-room building envelope + internal loads
}
\end{lstlisting}

\section{Data Flow Diagram}

\begin{verbatim}
  UI Input
     |
     v
  Redux Dispatch (action)
     |
     v
  Slice Reducer (state update)
     |
     v
  selectRdsData (createSelector, memoised)
     |
     +-- calculateSeasonLoad()    [seasonalLoads.ts]
     +-- calculateAirQuantities() [airQuantities.ts]
     +-- calculateAllSeasonOALoads() [outdoorAirLoad.ts]
     +-- calculateRequiredADP()   [psychro.ts]
     +-- calculateHeatingHumid()  [heatingHumid.ts]
     +-- calculatePipeSizing()    [pipeSizing.ts]
     +-- calculateAllSeasonStatePoints() [psychroStatePoints.ts]
     |
     v
  RDS Row Object (one per room)
     |
     v
  React Component Tree (re-render)
\end{verbatim}

\section{Directory Structure}

\begin{lstlisting}[language={}, caption={Source Directory Layout}]
client/src/
  app/
    store.ts           -- Redux store configuration
    rootReducer.ts     -- Slice composition
  components/
    Layout/            -- Header, TabNav, RoomSidebar
    UI/                -- InputField, InputGroup, NumberControl, StatCard
  constants/
    ashrae.ts          -- ASHRAE constants & conversion factors
    ashraeTables.ts    -- CLTD/CLF lookup tables, U-value presets
    isoCleanroom.ts    -- ISO 14644 classification data
    ventilation.ts     -- ASHRAE 62.1-2022 VRP implementation
  features/
    project/projectSlice.ts   -- Project, ambient, system design state
    ahu/ahuSlice.ts           -- AHU definitions
    room/roomSlice.ts         -- Room list & geometry
    room/roomActions.js       -- Cross-slice room thunks
    room/AhuAssignment.jsx    -- AHU-to-room assignment UI
    climate/climateSlice.ts   -- Seasonal outdoor conditions
    envelope/envelopeSlice.ts -- Building elements + internal loads
    envelope/BuildingShell.tsx -- Envelope editor component
    results/
      rdsSelector.ts          -- Master RDS orchestrator
      seasonalLoads.ts        -- Per-season load calculation
      airQuantities.ts        -- All CFM quantities
      outdoorAirLoad.ts       -- OA coil loads
      heatingHumid.ts         -- Heating + humidification
      pipeSizing.ts           -- CHW/HW pipe diameter selection
      psychroStatePoints.ts   -- AHU air stream state points
  hooks/
    useActiveRoom.js          -- Active room selection hook
    useNumberControl.ts       -- Numeric input stepper hook
    useProjectTotals.js       -- Project-wide aggregation hook
    useRdsRow.js              -- Single RDS row data hook
    useRoomSidebar.ts         -- Sidebar navigation hook
  pages/
    Home.jsx                  -- Landing page
    ProjectDetails.jsx        -- Project metadata entry
    AHUConfig.jsx             -- AHU equipment configuration
    RoomConfig.jsx            -- Room geometry & classification
    ClimateConfig.jsx         -- Outdoor design conditions
    EnvelopeConfig.jsx        -- Building envelope editor
    RDSPage.jsx               -- Room Data Sheet spreadsheet view
    ResultsPage.jsx           -- Project results & export
    rds/                      -- RDS sub-components
  utils/
    psychro.ts                -- Psychrometric calculations
    units.ts                  -- Unit conversions
    types.ts                  -- TypeScript interfaces
    envelopeCalc.ts           -- CLTD/CLF envelope calculations
    envelopeAggregator.ts     -- Envelope gain summation
    envelopeHelpers.ts        -- Envelope utility functions
    glazingCalc.ts            -- Glass & skylight solar gains
    isoValidation.ts          -- ISO/GMP compliance validation
    psychroValidation.ts      -- Psychrometric consistency checks
\end{lstlisting}


% ══════════════════════════════════════════════════════════════════════════════
\chapter{Entry Points \& Application Shell}
% ══════════════════════════════════════════════════════════════════════════════

\section{\file{client/src/main.jsx} --- Application Bootstrap}

The application entry point. Renders the React root with:
\begin{itemize}
  \item \code{StrictMode} --- Enables development-time double-rendering for side-effect detection.
  \item \code{Provider store=\{store\}} --- Makes the Redux store available to all descendant components via \code{useSelector} / \code{useDispatch}.
  \item \code{App} --- The root application component (routing).
\end{itemize}

\section{\file{client/src/App.jsx} --- Routing \& Layout}

Defines the application's route structure using React Router v7:

\begin{table}[H]
\centering
\caption{Application Routes}
\begin{tabularx}{\textwidth}{llX}
\toprule
\textbf{Path} & \textbf{Component} & \textbf{Description} \\
\midrule
\texttt{/}         & \code{Home}            & Public landing page (no header/nav) \\
\texttt{/project}  & \code{ProjectDetails}  & Project metadata and system design \\
\texttt{/rds}      & \code{RDSPage}         & Room Data Sheet spreadsheet \\
\texttt{/ahu}      & \code{AHUConfig}       & AHU equipment configuration \\
\texttt{/room}     & \code{RoomConfig}      & Room geometry and classification \\
\texttt{/climate}  & \code{ClimateConfig}   & Outdoor design conditions \\
\texttt{/envelope} & \code{EnvelopeConfig}  & Building envelope editor \\
\texttt{/results}  & \code{ResultsPage}     & Results summary and export \\
\texttt{*}         & Redirect to \texttt{/} & Catch-all \\
\bottomrule
\end{tabularx}
\end{table}

The \code{AppLayout} wrapper provides the fixed-viewport layout: \code{Header} + \code{TabNav} at natural height, \code{<main>} fills remaining viewport via \code{flex-1 overflow-auto min-h-0}.

\section{\file{client/src/app/store.ts} --- Redux Store}

Configures the Redux store with:
\begin{itemize}
  \item Root reducer composed from five slices.
  \item Dev tools enabled in non-production mode.
  \item Serializable check with \code{climate} slice ignored (transient \code{NaN} during numeric input).
  \item Automatic type inference for \code{RootState} and \code{AppDispatch}.
  \item Global \code{window.\_\_store} binding in development for console debugging.
\end{itemize}


% ══════════════════════════════════════════════════════════════════════════════
\chapter{Redux State Slices}
% ══════════════════════════════════════════════════════════════════════════════

\section{\file{features/project/projectSlice.ts} --- Project State}

\subsection{State Shape}

\begin{lstlisting}[language={}, caption={Project State Shape}]
state.project = {
  info: {
    projectName, projectLocation, customerName,
    consultantName, industry, keyAccountManager
  },
  ambient: {
    elevation,        // ft (0 = sea level)
    latitude,         // degrees (-90..90); default 28 (Delhi)
    dailyRange,       // F daily DB swing
    dryBulbTemp,      // C (reference only, NOT used in calcs)
    wetBulbTemp,      // C (reference only)
    relativeHumidity  // % (reference only)
  },
  systemDesign: {
    safetyFactor,           // % (default 10)
    ductHeatGain,           // % (default 5)
    bypassFactor,           // dimensionless (default 0.10)
    adp,                    // F (default 55)
    adpMode,                // 'manual' | 'calculated'
    fanHeat,                // % (default 5)
    returnFanHeat,          // % (default 5)
    humidificationTarget    // %RH (default 45)
  }
}
\end{lstlisting}

\subsection{Reducers}

\begin{itemize}
  \item \code{updateProjectInfo(\{field, value\})} --- Updates project metadata string fields with \code{trim()}.
  \item \code{updateAmbient(\{field, value\})} --- Parses, clamps to \code{AMBIENT\_BOUNDS}, and stores numeric ambient values.
  \item \code{updateSystemDesign(\{field, value\})} --- String fields (e.g.\ \code{adpMode}) bypass \code{parseFloat}; numeric fields are clamped to \code{SYSTEM\_DESIGN\_BOUNDS}.
  \item \code{resetProject()} --- Returns entire slice to \code{initialState}.
\end{itemize}

\subsection{Selectors}

14 exported selectors provide targeted access: \code{selectProjectInfo}, \code{selectAmbient}, \code{selectSystemDesign}, \code{selectElevation}, \code{selectLatitude}, \code{selectDailyRange}, \code{selectSafetyFactor}, \code{selectDuctHeatGain}, \code{selectBypassFactor}, \code{selectAdp}, \code{selectFanHeat}, \code{selectReturnFanHeat}, \code{selectHumidTarget}, \code{selectIndustry}.

\subsection{Key Design Decisions}

\begin{itemize}
  \item \textbf{ADP Scope:} Project ADP is the default; per-AHU overrides live in \code{ahuSlice}. Priority chain: AHU calculated $>$ AHU manual $>$ project default $>$ 55\textdegree F fallback.
  \item \textbf{Duct Heat Gain:} Applied \emph{additively} with safety factor per Excel row~80 method: $\text{mult} = 1 + (\text{safety} + \text{duct}) / 100$.
  \item \textbf{Ambient fields} (\code{dryBulbTemp}, \code{wetBulbTemp}, \code{relativeHumidity}) are brief reference values in \textdegree C --- they are \emph{not} used in load calculations. Seasonal design conditions live in \code{climateSlice}.
\end{itemize}

% ──────────────────────────────────────────────────────────────────────────────
\section{\file{features/ahu/ahuSlice.ts} --- AHU State}

Manages the list of Air Handling Units. Each AHU has identity fields (\code{id}, \code{name}, \code{tag}), a \code{type} (Recirculating, DOAS, MAU, FCU), capacity and airflow fields, and psychrometric override fields (\code{adp}, \code{adpMode}, \code{bypassFactor}).

\subsection{AHU Types}
\begin{description}
  \item[\code{Recirculating}] Standard recirculating AHU with partial outdoor air.
  \item[\code{DOAS}] Dedicated Outdoor Air System---100\% OA always.
  \item[\code{MAU}] Make-up Air Unit---100\% OA, typically for laboratories.
  \item[\code{FCU}] Fan Coil Unit---recirculating, no OA path.
\end{description}

The logic layer distinguishes only DOAS vs non-DOAS (\code{isDOAS = ahuType === 'DOAS'}).

\subsection{Reducers}
\code{addAHU}, \code{updateAHU}, \code{deleteAHU}. \warn{\code{deleteAHU} must never be called directly from UI code---use \code{deleteAhuWithCleanup} thunk from \file{roomActions.js}.}

% ──────────────────────────────────────────────────────────────────────────────
\section{\file{features/room/roomSlice.ts} --- Room State}

\subsection{State Shape}
\begin{lstlisting}[language={}, caption={Room State Shape}]
state.room = {
  activeRoomId: string | null,
  list: Room[]  // see Room interface in types.ts
}
\end{lstlisting}

Each \code{Room} contains: identity fields, geometry in SI (m\textsuperscript{2}, m\textsuperscript{3}), design targets (\code{designTemp} in \textdegree C, \code{designRH} in \%, \code{pressure} in Pa), ISO classification, ventilation category, ACPH constraints, exhaust air breakdown (general, BIBO, machine in CFM), and AHU assignment.

\subsection{Key Behaviours}
\begin{itemize}
  \item \code{designDB} (\textdegree F) is derived from \code{designTemp} (\textdegree C) and kept in sync by \code{updateRoom}.
  \item \code{classInOp} changes auto-update \code{minAcph} and \code{designAcph} from \code{ACPH\_RANGES}.
  \item Geometry fields (\code{length}, \code{width}, \code{height}) auto-derive \code{floorArea} and \code{volume}.
  \item \code{designRH = 0} is valid (battery dry rooms)---the \code{!= null} guard preserves zero.
  \item \code{deleteRoom} never removes the last room.
\end{itemize}

% ──────────────────────────────────────────────────────────────────────────────
\section{\file{features/climate/climateSlice.ts} --- Climate State}

Stores outdoor design conditions for three seasons. All temperatures in \textdegree F.

\begin{lstlisting}[language={}, caption={Climate State Shape}]
state.climate.outside = {
  summer:  { db: 109.9, rh: 19, time, month, gr, dp, wb },
  monsoon: { db: 95,    rh: 70, ... },
  winter:  { db: 45,    rh: 60, ... }
}
\end{lstlisting}

Fields \code{gr}, \code{dp}, \code{wb} are display-only derived fields computed at sea level (elevation = 0). The calculation layer recomputes these with actual site elevation.

Summer default: 109.9\textdegree F (43.3\textdegree C)---ASHRAE HOF 2021 Ch.14 Table~1, 0.4\% design dry-bulb for Delhi (28\textdegree N), correct tier for 24/7 critical facilities.

% ──────────────────────────────────────────────────────────────────────────────
\section{\file{features/envelope/envelopeSlice.ts} --- Envelope State}

Per-room envelope data stored in \code{state.envelope.byRoomId[roomId]}.

\subsection{Structure}
\begin{lstlisting}[language={}, caption={Room Envelope Structure}]
RoomEnvelope = {
  elements: {
    walls: [], roofs: [], glass: [],
    skylights: [], partitions: [], floors: []
  },
  internalLoads: {
    people:    { count, sensiblePerPerson, latentPerPerson },
    lights:    { wattsPerSqFt, useSchedule, ballastFactor },
    equipment: { kw, sensiblePct, latentPct, diversityFactor }
  },
  infiltration: {
    method: 'ach', achValue: 0, cfmValue: 0, doors: []
  }
}
\end{lstlisting}

\code{achValue} defaults to 0 (positively pressurised ISO/GMP rooms have zero infiltration by definition).

\subsection{Reducers}
\code{initializeRoom} (ISO-aware, guards against re-initialisation), \code{addEnvelopeElement}, \code{updateEnvelopeElement}, \code{removeEnvelopeElement}, \code{updateInternalLoad} (merge-update), \code{updateInfiltration}, \code{removeRoomEnvelope}.


% ══════════════════════════════════════════════════════════════════════════════
\chapter{Cross-Slice Thunks}
% ══════════════════════════════════════════════════════════════════════════════

\section{\file{features/room/roomActions.js}}

Three thunks for atomic cross-slice operations:

\begin{description}
  \item[\code{addNewRoom()}] Creates a room in \code{roomSlice} and initialises its envelope in \code{envelopeSlice} atomically. Uses \code{DEFAULT\_NEW\_ROOM\_CLASS = 'ISO 8'} as the single source of truth for both classification and ACPH defaults.

  \item[\code{deleteRoomWithCleanup(roomId)}] Removes room from \code{roomSlice} and its envelope from \code{envelopeSlice}.

  \item[\code{deleteAhuWithCleanup(ahuId)}] Clears all room assignments (\code{setRoomAhu(\{ahuId: null\})}) \emph{before} removing the AHU from \code{ahuSlice}. Order matters: clear first, then remove.
\end{description}


% ══════════════════════════════════════════════════════════════════════════════
\chapter{Constants \& Reference Data}
% ══════════════════════════════════════════════════════════════════════════════

\section{\file{constants/ashrae.ts} --- ASHRAE Constants}

Central repository of \ashrae{} Handbook constants. Exported as a frozen \code{as const} object. Key groups:

\begin{table}[H]
\centering
\caption{Key ASHRAE Constants}
\begin{tabularx}{\textwidth}{lrX}
\toprule
\textbf{Constant} & \textbf{Value} & \textbf{Source / Notes} \\
\midrule
\code{SENSIBLE\_FACTOR\_SEA\_LEVEL} & 1.08 & $Q_s = 1.08 \times \text{CFM} \times \Delta T$ \\
\code{LATENT\_FACTOR\_SEA\_LEVEL}   & 0.68 & $Q_l = 0.68 \times \text{CFM} \times \Delta\text{gr/lb}$ \\
\code{LATENT\_FACTOR\_LB}           & 4775 & $= 0.075 \times 60 \times 1061$ \\
\code{LATENT\_HFG\_BTU\_LB}         & 1061 & $h_{fg}$ at 32\textdegree F reference \\
\code{BTU\_PER\_TON}                & 12000 & 1 TR = 12,000 BTU/hr \\
\code{M2\_TO\_FT2}                  & 10.7639 & NIST exact \\
\code{PROCESS\_SAFETY\_FACTOR}      & 1.25 & GMP Annex 1:2022 \S4.23 \\
\code{PROCESS\_DIVERSITY\_FACTOR}   & 0.75 & Fallback only \\
\bottomrule
\end{tabularx}
\end{table}

\warn{Never use \code{SENSIBLE\_FACTOR\_SEA\_LEVEL} or \code{LATENT\_FACTOR\_SEA\_LEVEL} directly in calculations. Use \code{sensibleFactor(elev)} / \code{latentFactor(elev)} from \file{psychro.ts}.}

\section{\file{constants/isoCleanroom.ts} --- ISO 14644 Classification}

Contains:
\begin{itemize}
  \item \code{ISO\_PARTICLE\_LIMITS} --- ISO 14644-1:2015 particle concentration limits for ISO~1--9.
  \item \code{ACPH\_RANGES} --- IEST-RP-CC012.2 air change rates (min, design, max, OA fraction, flow type).
  \item \code{ISO\_CLASS\_DATA} --- Full classification data including GMP grade mapping, typical use, pressure, temperature, and humidity ranges.
  \item \code{GMP\_GRADE\_MAPPING} --- GMP Annex 1:2022 Grade~A--D with ISO at-rest/in-operation requirements and mandated ACPH.
\end{itemize}

\section{\file{constants/ventilation.ts} --- ASHRAE 62.1-2022 VRP}

Implements the Ventilation Rate Procedure with 12 industry-specific categories:

\begin{table}[H]
\centering
\caption{Ventilation Categories (Selected)}
\begin{tabularx}{\textwidth}{llrrX}
\toprule
\textbf{Key} & \textbf{Label} & $R_p$ & $R_a$ & \textbf{Min ACH} \\
\midrule
\code{general}           & General / Office           & 5    & 0.06 & 0 \\
\code{pharma}            & Pharma Cleanroom           & 5    & 0.18 & 20 \\
\code{semicon}           & Semiconductor Fab          & 5    & 0.18 & 6 \\
\code{battery}           & Battery --- Li-ion         & 5    & 0.18 & 10 \\
\code{battery-leadacid}  & Battery --- Lead-Acid      & 5    & 1.00 & 12 \\
\code{gowning}           & Gowning Room               & 5    & 0.10 & 10 \\
\code{canteen}           & Canteen / Cafeteria        & 7.5  & 0.18 & 0 \\
\bottomrule
\end{tabularx}
\end{table}

Key exports: \code{calculateVbz()} (zone OA: $V_{bz} = R_p \times P_z + R_a \times A_z$) and \code{calculateMinAchCfm()} (regulatory ACH floor).


% ══════════════════════════════════════════════════════════════════════════════
\chapter{Psychrometric Engine}
% ══════════════════════════════════════════════════════════════════════════════

\section{\file{utils/psychro.ts} --- Core Psychrometric Functions}

The psychrometric engine uses \textbf{ASHRAE Hyland-Wexler} saturation pressure equations (HOF 2021 Ch.1, Eq.~3 for ice, Eq.~5 for liquid water), replacing the less accurate Magnus approximation used in earlier versions.

\subsection{Saturation Pressure}

\begin{equation}
  \ln(P_{ws}) = \frac{C_1}{T} + C_2 + C_3 T + C_4 T^2 + C_5 T^3 + C_6 T^4 + C_7 \ln(T) \quad \text{(ice, } T < 273.15\text{ K)}
\end{equation}

\begin{equation}
  \ln(P_{ws}) = \frac{C_8}{T} + C_9 + C_{10} T + C_{11} T^2 + C_{12} T^3 + C_{13} \ln(T) \quad \text{(liquid, } T \geq 273.15\text{ K)}
\end{equation}

\subsection{Public API}

\begin{longtable}{p{6cm}p{8cm}}
\toprule
\textbf{Function} & \textbf{Description} \\
\midrule
\endhead
\code{altitudeCorrectionFactor(elevFt)} & $C_f = (1 - 6.8754 \times 10^{-6} \times \text{elev})^{5.2559}$ \\
\code{sitePressure(elevFt)} & $P_{\text{atm}} = 1013.25 \times C_f$ (hPa) \\
\code{sensibleFactor(elevFt)} & $C_s = 1.08 \times C_f$ \\
\code{latentFactor(elevFt)} & $C_l = 0.68 \times C_f$ \\
\code{calculateGrains(dbF, rh, elevFt)} & $W = 0.62198 \times E / (P_{\text{atm}} - E)$ (gr/lb) \\
\code{calculateRH(dbF, grains, elevFt)} & Reverse of \code{calculateGrains} \\
\code{calculateDewPoint(dbF, rh)} & Bisection on H-W curve; returns null for $\text{RH} \leq 0$ \\
\code{grainsFromDewPoint(dpF, elevFt)} & $W$ from dew/frost point temperature \\
\code{calculateEnthalpy(dbF, grains)} & $h = 0.240t + W(h_{fg_0} + 0.444t)$ \\
\code{calculateWetBulb(dbF, rh, elevFt)} & Bisection on ASHRAE Eq.~35 ($\pm 0.01$\textdegree C) \\
\code{calculateSpecificVolume(dbF, gr, elevFt)} & $v = 0.370486 \times T_R / P_{\text{psia}} \times (1 + 1.608W)$ \\
\code{calculateAdpFromLoads(...)} & Back-calculates ADP from room sensible load \\
\code{calculateRequiredADP(...)} & ESHF line / saturation curve intersection \\
\bottomrule
\end{longtable}

\subsection{Required ADP --- ESHF Analysis}

Three possible outcomes:
\begin{description}
  \item[\code{found}] Bisection succeeded. \code{requiredADP} is the coil surface temperature needed. Compare with plant ADP.
  \item[\code{sensible\_only}] Required ADP $< 32$\textdegree F or ESHF $\geq 99.5\%$. Any standard CHW coil works.
  \item[\code{no\_solution}] ESHF line does not intersect saturation curve. Supplemental dehumidification required.
\end{description}

\section{\file{utils/units.ts} --- Unit Conversion Library}

Pure, zero-dependency conversion library. Organised by family: area, volume, temperature, energy/power, humidity, airflow, pressure, mass flow, and length.

Key design decisions:
\begin{itemize}
  \item Temperature conversions use \code{numOrNull} guard (returns \code{null} for invalid input) because 0\textdegree F is a valid temperature.
  \item Positive-value conversions (area, airflow) use \code{num()} guard (returns 0 for invalid input).
  \item Includes \code{PSYCHRO\_SANITY\_CHECKS} reference values for regression testing.
\end{itemize}


% ══════════════════════════════════════════════════════════════════════════════
\chapter{Calculation Modules}
% ══════════════════════════════════════════════════════════════════════════════

\section{\file{features/results/seasonalLoads.ts} --- Seasonal Load Calculation}

Computes per-room, per-season sensible and latent loads (BTU/hr).

\subsection{Load Components}

\textbf{Sensible loads:}
\begin{enumerate}
  \item \textbf{Envelope} --- CLTD/CLF method via \code{envelopeAggregator.ts}
  \item \textbf{People} --- $\text{sensible/person} \times \text{count}$ (HOF Ch.18 Table~1)
  \item \textbf{Lighting} --- $\text{W/ft}^2 \times \text{area} \times 3.412 \times \text{schedule} \times \text{ballast}$
  \item \textbf{Equipment} --- $\text{kW} \times 3412.14 \times \text{sensible\%} \times \text{diversity}$
  \item \textbf{Infiltration} --- $C_s \times \text{CFM} \times \Delta T$
\end{enumerate}

\textbf{Latent loads:}
\begin{enumerate}
  \item People latent heat
  \item Equipment latent fraction
  \item Infiltration --- $C_l \times \text{CFM} \times \Delta\text{gr/lb}$
\end{enumerate}

\subsection{Safety Factor Policy}

\begin{equation}
  \text{ERSH} = \text{rawSensible} \times \underbrace{\left(1 + \frac{\text{safety} + \text{ductGain}}{100}\right)}_{\text{safetyMult}} \times \underbrace{1.25}_{\text{GMP rooms only}}
\end{equation}

\begin{equation}
  \text{ERLH} = \text{rawLatent} \quad \text{(no safety factor on latent)}
\end{equation}

% ──────────────────────────────────────────────────────────────────────────────
\section{\file{features/results/airQuantities.ts} --- Airflow Calculations}

\subsection{Supply Air Hierarchy}

\begin{equation}
  \text{supplyAir} = \max(\text{thermalCFM}, \text{designACPH\_CFM}, \text{regulatoryACPH\_CFM}, \text{minACPH\_CFM})
\end{equation}

Where:
\begin{align}
  \text{thermalCFM} &= \frac{\text{ERSH}}{C_s \times (1 - BF) \times (T_{\text{room}} - T_{\text{ADP}})} \\
  \text{designACPH\_CFM} &= V_{\text{ft}^3} \times \text{designACPH} / 60 \\
  \text{regulatoryACPH\_CFM} &= \code{calculateMinAchCfm}(\text{ventCategory}, V_{\text{ft}^3})
\end{align}

\subsection{Fresh Air (ASHRAE 62.1 VRP)}

\begin{equation}
  V_{bz} = R_p \times P_z + R_a \times A_z
\end{equation}

Exhaust compensation: $\text{freshAir} = \max(V_{bz}, \text{totalExhaust})$.

For DOAS units: $\text{freshAir} = \text{supplyAir}$ (100\% OA).

\subsection{Mass Balance}

\begin{equation}
  \text{returnAir} = \max(0, \text{supplyAir} - \text{freshAirCheck})
\end{equation}

% ──────────────────────────────────────────────────────────────────────────────
\section{\file{features/results/rdsSelector.ts} --- RDS Orchestrator}

The master memoised selector (\code{createSelector}) that assembles the complete RDS row for every room. Pure orchestrator---delegates all calculations:

\begin{enumerate}
  \item \textbf{Step 1:} Seasonal loads (3 seasons $\times$ each room)
  \item \textbf{Step 2:} Air quantities (supply, fresh, exhaust, return)
  \item \textbf{Step 3:} Outdoor air coil loads (enthalpy method)
  \item \textbf{Step 4:} Peak cooling season selection + grand total
  \item \textbf{Step 4b:} Required ADP from ESHF line analysis
  \item \textbf{Step 4c:} Reheater requirement (ASHRAE Ch.18 \S17.3)
  \item \textbf{Step 5:} Heating + humidification sizing
  \item \textbf{Step 6:} CHW/HW pipe sizing
  \item \textbf{Step 7:} Psychrometric state points
  \item \textbf{Step 8:} Derived seasonal fields (pickup, terminal heating)
\end{enumerate}

\subsection{Reheater Logic}

When $\text{roomESHF} < \text{minESHF}$, a reheater is required:

\begin{equation}
  \text{minESHF} = \frac{C_s \times DR}{C_s \times DR + C_l \times \max(0, W_{\text{room}} - W_{\text{sat}}(T_{\text{ADP}}))}
\end{equation}

\begin{equation}
  \text{reheatBTU} = \frac{\text{minESHF} \times \text{ERTH}_{\text{safe}} - \text{ERSH}_{\text{safe}}}{1 - \text{minESHF}}
\end{equation}

\subsection{ADP Resolution Priority}

\begin{enumerate}
  \item AHU mode = \code{calculated} $\rightarrow$ \code{calculateAdpFromLoads()}
  \item \code{ahu.adp > 0} $\rightarrow$ per-AHU manual override
  \item \code{systemDesign.adp} $\rightarrow$ project default
  \item 55\textdegree F hardcoded fallback
\end{enumerate}


% ══════════════════════════════════════════════════════════════════════════════
\chapter{Envelope Calculation Modules}
% ══════════════════════════════════════════════════════════════════════════════

\section{\file{utils/envelopeCalc.ts}}

CLTD/CLF-based heat gain calculations for opaque building elements:
\begin{itemize}
  \item \code{calcWallGain(element, climate, tRoom, season)} --- Wall conduction with CLTD, latitude/month correction, and daily range correction.
  \item \code{calcRoofGain(element, climate, tRoom, season)} --- Roof conduction via flat-roof CLTD tables.
  \item \code{calcPartitionGain(element, tRoom, season)} --- Steady-state conduction $Q = U \times A \times \Delta T$ with season-dependent adjacent temperatures.
  \item \code{calcInfiltrationGain(inf, room, volFt3, ...)} --- Sensible and latent infiltration loads.
\end{itemize}

\section{\file{utils/glazingCalc.ts}}

Solar heat gain through glass and skylights:
\begin{itemize}
  \item Conduction: $Q_c = U \times A \times \text{CLTD}_{\text{glass}}$
  \item Solar: $Q_s = A \times \text{SHGF} \times \text{SHGC} \times \text{CLF}$
  \item Total: $Q = Q_c + Q_s$
\end{itemize}

SHGC preferred over SC ($\text{SHGC} = 0.87 \times \text{SC}$).

\section{\file{utils/envelopeAggregator.ts}}

Sums gains across all element categories for one room in one season. Calls individual \code{calcXxxGain()} functions from \file{envelopeCalc.ts} and \file{glazingCalc.ts}.

\section{\file{features/envelope/BuildingShell.tsx}}

Interactive tabbed UI for editing building envelope elements. Six categories (walls, roofs, glass, skylights, partitions, floors) with per-element ASHRAE preset selection, U-value input, area input, and per-season BTU/hr gain badges.


% ══════════════════════════════════════════════════════════════════════════════
\chapter{Utility Modules}
% ══════════════════════════════════════════════════════════════════════════════

\section{\file{utils/types.ts} --- TypeScript Interfaces}

Defines all shared interfaces:

\begin{table}[H]
\centering
\caption{Key TypeScript Interfaces}
\begin{tabularx}{\textwidth}{lX}
\toprule
\textbf{Interface} & \textbf{Purpose} \\
\midrule
\code{Room}              & Room identity, geometry, design targets, classification \\
\code{RoomState}         & \code{\{activeRoomId, list: Room[]\}} \\
\code{AHU}               & AHU equipment specification \\
\code{AhuState}          & \code{\{list: AHU[]\}} \\
\code{ProjectInfo}       & Project metadata \\
\code{AmbientParams}     & Site reference conditions \\
\code{SystemDesign}      & HVAC system sizing parameters \\
\code{ProjectState}      & \code{\{info, ambient, systemDesign\}} \\
\code{EnvelopeElement}   & Building element with U-value, area, orientation \\
\code{InternalPeople}    & Occupant heat gain parameters \\
\code{InternalLights}    & Lighting heat gain parameters \\
\code{InternalEquipment} & Equipment heat gain parameters \\
\code{RoomInfiltration}  & Infiltration method and values \\
\code{RoomEnvelope}      & Elements + internal loads + infiltration \\
\code{EnvelopeState}     & \code{\{byRoomId: Record<string, RoomEnvelope>\}} \\
\code{SeasonCondition}   & Outdoor condition for one season \\
\code{ClimateState}      & Three-season outdoor conditions \\
\bottomrule
\end{tabularx}
\end{table}

\code{RootState} is re-exported from \code{store.ts} via \code{ReturnType<typeof store.getState>} for automatic type inference.

\section{\file{utils/isoValidation.ts}}

ISO 14644 and GMP compliance validation. Checks room classification consistency, ACPH sufficiency, pressure cascade compliance, and GMP grade mapping validity.

\section{\file{utils/psychroValidation.ts}}

Psychrometric consistency checks. Validates that room humidity targets are achievable given the selected ADP and bypass factor.


% ══════════════════════════════════════════════════════════════════════════════
\chapter{UI Components}
% ══════════════════════════════════════════════════════════════════════════════

\section{Layout Components}

\begin{description}
  \item[\file{components/Layout/Header.tsx}] Application header bar with project branding.
  \item[\file{components/Layout/HeaderBadge.tsx}] ASHRAE standard compliance badge.
  \item[\file{components/Layout/TabNav.tsx}] Horizontal tab navigation (Project, RDS, AHU, Room, Climate, Envelope, Results).
  \item[\file{components/Layout/RoomSidebar.tsx}] Vertical room list sidebar for multi-room navigation.
  \item[\file{components/Layout/RoomSidebarItem.tsx}] Individual room entry in sidebar.
\end{description}

\section{UI Primitives}

\begin{description}
  \item[\file{components/UI/InputField.tsx}] Labelled text/number input with validation state.
  \item[\file{components/UI/InputGroup.tsx}] Grouped inputs with shared label.
  \item[\file{components/UI/NumberControl.tsx}] Numeric stepper with increment/decrement buttons.
  \item[\file{components/UI/StatCard.tsx}] Display card for computed statistics.
\end{description}

\section{\file{features/room/AhuAssignment.jsx}}

Radio-button single-select AHU assignment component. Replaces the previous multi-select checkbox UI that was architecturally inconsistent---all downstream consumers read only \code{assignedAhuIds[0]}.


% ══════════════════════════════════════════════════════════════════════════════
\chapter{Custom Hooks}
% ══════════════════════════════════════════════════════════════════════════════

\begin{description}
  \item[\file{hooks/useActiveRoom.js}] Returns the currently selected room object and dispatch helpers for room field updates.

  \item[\file{hooks/useNumberControl.ts}] Manages numeric input state with increment/decrement, min/max bounds, and step size.

  \item[\file{hooks/useProjectTotals.js}] Aggregates RDS data across all rooms into project-level totals (total TR, total CFM, per-AHU summaries).

  \item[\file{hooks/useRdsRow.js}] Returns the full computed RDS row object for a single room, including dispatch helpers for inline editing.

  \item[\file{hooks/useRoomSidebar.ts}] Manages room sidebar state (active room, add/delete room, room ordering).
\end{description}


% ══════════════════════════════════════════════════════════════════════════════
\chapter{Pages}
% ══════════════════════════════════════════════════════════════════════════════

\section{\file{pages/Home.jsx} --- Landing Page}
Public-facing landing page without the application header/nav. Entry point for new projects.

\section{\file{pages/ProjectDetails.jsx} --- Project Configuration}
Three sections: project metadata (name, location, customer), site ambient conditions (elevation, latitude, daily range), and system design parameters (safety factor, BF, ADP, fan heat).

\section{\file{pages/AHUConfig.jsx} --- AHU Equipment}
Add, edit, and delete AHUs. Fields: type (Recirculating/DOAS/MAU/FCU), capacity, airflow, psychrometric overrides, filter class, location, notes.

\section{\file{pages/RoomConfig.jsx} --- Room Setup}
Left column: room identity, geometry, design targets, ISO classification, ventilation category, ACPH settings, exhaust air breakdown. Right column: AHU assignment.

\section{\file{pages/ClimateConfig.jsx} --- Outdoor Conditions}
Three-season (summer, monsoon, winter) outdoor design conditions. Each season: DB, RH, time of peak, month. Derived fields (gr, dp, wb) shown as read-only.

\section{\file{pages/EnvelopeConfig.jsx} --- Building Envelope}
Hosts the \code{BuildingShell} component. Six categories of envelope elements with ASHRAE construction presets and per-season BTU/hr gain display.

\section{\file{pages/RDSPage.jsx} --- Room Data Sheet}
Spreadsheet-style view of all computed RDS data. Groups rooms by AHU. Sections: setup, loads, psychrometric state points, results.

\section{\file{pages/ResultsPage.jsx} --- Project Results}
Project-level aggregation: total cooling capacity (TR), total heating capacity (kW), per-AHU summaries, check figure (ft\textsuperscript{2}/TR), and Excel export.


% ══════════════════════════════════════════════════════════════════════════════
\chapter{Deployment}
% ══════════════════════════════════════════════════════════════════════════════

\section{Build \& Deploy}

\begin{lstlisting}[language=bash, caption={Build Commands}]
cd client
npm install
npm run build    # Production bundle
npm run dev      # Development server with HMR
npm run test     # Vitest test suite
\end{lstlisting}

\section{Vercel Configuration}

SPA rewrite rule in \file{vercel.json}:
\begin{lstlisting}[language={}, caption={Vercel Rewrite Rule}]
{
  "rewrites": [{ "source": "/(.*)", "destination": "/index.html" }]
}
\end{lstlisting}

All client-side routes are handled by React Router; the server always serves \code{index.html}.


% ══════════════════════════════════════════════════════════════════════════════
\chapter{Engineering Standards Reference}
% ══════════════════════════════════════════════════════════════════════════════

\begin{longtable}{p{5cm}p{9cm}}
\toprule
\textbf{Standard} & \textbf{Usage in Application} \\
\midrule
\endhead
ASHRAE HOF 2021, Ch.1     & Psychrometric equations (Hyland-Wexler) \\
ASHRAE HOF 2021, Ch.18    & CLTD/CLF load calculation methodology \\
ASHRAE 62.1-2022          & Ventilation Rate Procedure (VRP) \\
ASHRAE 90.1-2022          & Energy standard (lighting power density) \\
ISO 14644-1:2015          & Cleanroom classification (particle limits) \\
ISO 14644-4:2022          & Cleanroom design \& construction \\
GMP Annex 1:2022          & Pharmaceutical HVAC requirements \\
IEST-RP-CC012.2           & Cleanroom air change rates \\
NFPA 855                  & Battery energy storage systems \\
IFC \S1206                & Fire code --- battery rooms \\
OSHA 29 CFR 1926.403(i)   & Lead-acid battery charging ventilation \\
IEEE 1184-2006            & Battery installation sizing \\
SEMI S2-0200              & Semiconductor facility safety \\
ISPE Baseline Guide Vol.5 & Pharmaceutical cleanroom HVAC \\
\bottomrule
\end{longtable}


% ══════════════════════════════════════════════════════════════════════════════
\chapter{Appendix: Key Equations}
% ══════════════════════════════════════════════════════════════════════════════

\section{Sensible Load}
\begin{equation}
  Q_s = C_s \times \text{CFM} \times \Delta T \quad [\text{BTU/hr}]
\end{equation}
where $C_s = 1.08 \times C_f$ and $C_f = (1 - 6.8754 \times 10^{-6} \times \text{elev}_{\text{ft}})^{5.2559}$.

\section{Latent Load}
\begin{equation}
  Q_l = C_l \times \text{CFM} \times \Delta W \quad [\text{BTU/hr}]
\end{equation}
where $C_l = 0.68 \times C_f$ and $\Delta W$ is in gr/lb.

\section{Humidity Ratio}
\begin{equation}
  W = \frac{0.62198 \times E}{P_{\text{atm}} - E} \quad [\text{kg/kg}]
\end{equation}
where $E = (\text{RH}/100) \times P_{ws}(T)$ and $P_{ws}$ is from Hyland-Wexler.

\section{Enthalpy}
\begin{equation}
  h = 0.240 \cdot t + W \cdot (1061 + 0.444 \cdot t) \quad [\text{BTU/lb dry air}]
\end{equation}

\section{Ventilation Zone Breathing Rate}
\begin{equation}
  V_{bz} = R_p \times P_z + R_a \times A_z \quad [\text{CFM}]
\end{equation}

\section{Thermal Supply Air}
\begin{equation}
  \text{CFM}_{\text{thermal}} = \frac{\text{ERSH}}{C_s \times (1 - BF) \times (T_{\text{room}} - T_{\text{ADP}})}
\end{equation}

\section{Effective Room Sensible Heat Factor}
\begin{equation}
  \text{ESHF} = \frac{Q_s^{\text{total}}}{Q_s^{\text{total}} + Q_l^{\text{total}}}
\end{equation}

\section{Reheater Capacity}
\begin{equation}
  Q_{\text{reheat}} = \frac{\text{minESHF} \times \text{ERTH} - \text{ERSH}}{1 - \text{minESHF}} \quad [\text{BTU/hr}]
\end{equation}

\section{Cooling Capacity}
\begin{equation}
  \text{TR} = \frac{\text{ERSH} + \text{ERLH} + Q_{\text{OA}} + Q_{\text{fan}} + Q_{\text{reheat}}}{12000}
\end{equation}

% ══════════════════════════════════════════════════════════════════════════════

\vfill
\begin{center}
\rule{0.5\textwidth}{0.4pt}\\[0.5em]
{\small\textcolor{codegray}{End of Documentation --- HeatLoad HVAC Calculation Tool v2.8}}
\end{center}

\end{document}
